% New complete positivity calculation answering the user's K comparable q idea.
\section{A strictly positive macroscopic return at the exact order $K=q+1$}
\label{sec:KPR}

\paragraph{Exploratory comparison status.}
This proof evaluates the auxiliary family of IKO, retaining its exact
relation to the original source. Its tested choice $K=q+1$ is an
externally selected comparison parameter. The original tensor source
and its canonical mixed-cohomology maps remain those of order $k$;
this calculation supplies no rule selecting a new canonical order.
The simultaneous baseline and return identities KPR15--16 are kept
in full, as required to use the comparison without an ad hoc insertion.

Retain the proved equilibrium $\rho_0(x)dx$ for the exact field
$V_0(x)=\pi\sqrt x-2\log x$, with support $[a_*,b_*]$ inside
$(0,\infty)$. Here $a_*,b_*$ denote these equilibrium endpoints
only; the earlier scalar error allowances retain their earlier
decorated notation. Write
\[
 \begin{gathered}
 M_1=\int\sqrt x\,\rho_0(x)dx=6/\pi,\qquad
 M_2=\int x\,\rho_0(x)dx,\\
 M_3=\int x^{3/2}\rho_0(x)dx,
 \quad\mathcal L=\int\log x\,\rho_0(x)dx.
 \end{gathered}
 \tag{KPR1}
\]
The first identity is EIQ31 with its original parameters. All moments
here are finite because the equilibrium support is compact and bounded
away from zero. We shall prove the sign of the actual independent-order
coefficient using these integrals and explicit finite inequalities.

\subsection{A further exact moment identity from a polynomial transport}
For sufficiently small real $t$ of either sign, the map
$T_t(x)=x+tx^2$ is positive and strictly increasing on $[a_*,b_*]$.
For example $|t|<1/(4b_*)$ ensures both assertions. Its probability
pushforward is therefore an admissible competitor in the exact
variational definition of $F_0$. For $x\ne y$ its pair difference is
\[
 T_t(x)-T_t(y)=(x-y)(1+t(x+y)),\qquad 1+t(x+y)>0.
\]
The change in pair energy is exactly the integral of
$\log(1+t(x+y))$, and its derivative at zero is $2M_2$.
Differentiation is justified by the uniform positive lower bound on
the last factor and the bounded support. The derivative of the
potential integral is
\[
 \int x^2V_0'(x)\rho_0(x)dx
       =\frac\pi2M_3-2M_2.
\]
The original equilibrium maximizes the functional, and zero is an
interior value of this explicit admissible parameter. Hence its
derivative is zero. We have proved
\[
 \boxed{M_3=\frac8\pi M_2.}
 \tag{KPR2}
\]
Applying Cauchy's inequality to the two functions $x^{1/4}$ and
$x^{3/4}$ in this original probability space gives
$M_2^2\leq M_1M_3$. Since $M_2>0$, KPR1--2 imply the finite bound
\[
 \boxed{0<M_2\leq\frac{48}{\pi^2}<\frac{16}{3}.}
 \tag{KPR3}
\]
The last inequality uses $\pi>3$. A direct rational verification is
\[
 \pi=4\int_0^1\frac{dx}{1+x^2}
 >4\sum_{j=0}^{7}\frac{(-1)^j}{2j+1}
 =\frac{135904}{45045}>3.
\]
The strict integral inequality follows by retaining the positive
remainder $x^{16}/(1+x^2)$ in the finite geometric identity.

\subsection{A specified dilated competitor for the full marked field}
For $\lambda\geq0$ retain exactly
\[
 h_\lambda(x)=\int_0^\lambda\log(x+4u^2)du,
 \quad V_\lambda(x)=\pi\sqrt x-2\log x-2h_\lambda(x),
 \quad F_\lambda=\sup_\eta\mathcal E_{V_\lambda}(\eta),
 \tag{KPR4}
\]
where $\mathcal E_V(\eta)=\iint\log|x-y|d\eta^2-\int Vd\eta$
has the precise single-kernel interpretation of MFG12.
For any $s>0$ take the explicit pushforward $x\mapsto s^2x$ of
$\rho_0dx$. Its value relative to $F_0$ is exactly
\[
 \begin{split}
 \mathcal E_{V_\lambda}((s^2\cdot)_*\rho_0)-F_0
 ={}&(6+4\lambda)\log s-6(s-1)+2\lambda\mathcal L\\
 &+2\int_0^\lambda\int
       \log\left(1+\frac{4u^2}{s^2x}\right)\rho_0(x)dx\,du.
 \end{split}
 \tag{KPR5}
\]
To verify every term, the pair energy gains $2\log s$;
the original $-\pi\sqrt x$ term gains $-\pi(s-1)M_1=-6(s-1)$;
the original $2\log x$ gains $4\log s$.
Finally split the full marked logarithm only by its exact identity
$\log(s^2x+4u^2)=2\log s+\log x+
\log(1+4u^2/(s^2x))$ and integrate it over its original interval.
All terms are finite on the specified support.

For $A\geq0$, the function $x\mapsto\log(1+A/x)$ is convex
on $(0,\infty)$: for $A>0$ its second derivative is
$A(2x+A)/(x^2(x+A)^2)>0$, and for $A=0$ it is zero.
Jensen's inequality followed by the elementary integral inequality
$\log(1+v)\geq v/(1+v)$ for $v\geq0$ therefore gives
\[
 \begin{split}
 2\int_0^\lambda\int\log\left(1+\frac{4u^2}{s^2x}\right)
                      \rho_0(x)dx\,du
 &\geq2\int_0^\lambda\log\left(1+\frac{4u^2}{s^2M_2}\right)du\\
 &\geq\frac{8\lambda^3}{3(s^2M_2+4\lambda^2)}.
 \end{split}
 \tag{KPR6}
\]
The last line retains the numerator $4u^2$, and bounds its denominator
$s^2M_2+4u^2$ from above by $s^2M_2+4\lambda^2$ before integration.
The logarithmic inequality follows, for example, by differentiating
$\log(1+v)-v/(1+v)$, whose derivative is $v/(1+v)^2\geq0$.

Use precisely $\lambda=1/4$ and $s=7/6$. Equations KPR3 and KPR6 give
\[
 2\int_0^{1/4}\int\log\left(1+\frac{4u^2}{(7/6)^2x}\right)
                          \rho_0(x)dx\,du
 \geq\frac{9}{1622},
 \tag{KPR7}
\]
since $(7/6)^2(16/3)+1/4=811/108$ and
$8(1/4)^3/[3(811/108)]=9/1622$.
The supremum $F_{1/4}$ is at least the value of this stated competitor.
Consequently
\[
 F_{1/4}-F_0
 \geq\frac12\mathcal L+7\log(7/6)-1+\frac9{1622}.
 \tag{KPR8}
\]

\subsection{A positive rational lower bound, with exact logarithmic errors}
The proved original recurrence bound and its exact limit GEL18a are
\[
 2\mathcal L-4(2\log2-1)=\mathcal C_\Gamma
      \geq4\log\frac{27}{8\pi}.
\]
Thus
$\mathcal L/2\geq\log(27/(2\pi))-1$
after combining the displayed logarithms. The exact polynomial
division identity
$22/7-\pi=\int_0^1x^4(1-x)^4/(1+x^2)dx>0$ implies
\[
 \frac12\mathcal L>\log(189/44)-1.
 \tag{KPR9}
\]
The literal source integral MFG20 at this same parameter is
\[
 Q(1/4)=\frac{25}{4}\log(5/2)-\frac94\log(3/2)
                      -4\log2-\frac32.
 \tag{KPR10}
\]
Combining KPR8--10 with the proved definition
$\Psi(1/4)=F_{1/4}-F_0-Q(1/4)$ yields
\[
 \Psi(1/4)>
 \log(189/44)+7\log(7/6)-\frac{25}{4}\log(5/2)
       +\frac94\log(3/2)+4\log2-\frac12+\frac9{1622}.
 \tag{KPR11}
\]

Here are finite rational bounds for every logarithm in this expression.
For $x>1$ set $r=(x-1)/(x+1)$. Integration of
$2/(1-t^2)$ from $0$ to $r$ gives
\[
 S_{20}(r)=2\sum_{j=0}^{19}\frac{r^{2j+1}}{2j+1}
 <\log x
 <S_{20}(r)+\frac{2r^{41}}{41(1-r^2)}.
 \tag{KPR12}
\]
The inequalities retain the positive remaining series and bound
each of its denominators below by $41$. For the five rational
arguments used here, substituting the displayed $r$ gives the
following exact rational enclosures. Every entry in the last two
columns is an integer, to be divided by $10^8$.
\[
 \begin{array}{c|c|r|r}
 x&r&\text{lower numerator}&\text{upper numerator}\\\hline
 189/44&145/233&145755738&145755739\\
 7/6&1/13&15415067&15415068\\
 5/2&3/7&91629073&91629074\\
 3/2&1/5&40546510&40546511\\
 2&1/3&69314718&69314719
 \end{array}
 \tag{KPR13}
\]
For completeness these rational comparisons require only
$n_-/10^8\leq S_{20}(r)$ and
$S_{20}(r)+2r^{41}/[41(1-r^2)]\leq n_+/10^8$ in each row.
Their denominators are positive; inserting its displayed rational
$r$ and multiplying the denominators verifies each finite inequality.
The accompanying exact-arithmetic script performs exactly these ten
stated integer comparisons and reproduces the following final
fraction; its proof inputs are the explicit sum and remainder KPR12.
Taking lower entries at positive coefficients and the upper entry
at the negative coefficient in KPR11 gives
\[
 \boxed{\Psi(1/4)>\frac{9279677}{40550000000}
                     >\frac1{5000}>0.}
 \tag{KPR14}
\]
Indeed the difference between the middle two fractions is
$1169677/40550000000>0$. This sign proof uses no numerical
quadrature of the equilibrium or finite matrix experiment.

\subsection{The original packet at the requested comparable order}
At $L=q/4$, $K=q+1$, the actual homogeneous-polynomial theorem MFG23
and the exact original-packet transport IKO17 prove
\[
 \boxed{\frac{\mathfrak T_{k;q/4}}{q^2}
           \longrightarrow\Psi(1/4)>\frac1{5000}.}
 \tag{KPR15}
\]
This is the positive return of the original packet with the new
independent Gamma order. Both original high ranks, the low relation
factor, and the entire $f_{q/4}$ occur in the finite formula IKO13.

The exact same-order identities IKO18--19, the proved original
baseline BSL16, and $\delta_k^\sigma/(kq)\to0$ now give
\[
 \boxed{\begin{aligned}
 \frac{\mathcal B_k^{(q+1)}}{q^2}&\longrightarrow C_B+\Psi(1/4),\\
 \frac{\mathcal B_k^{\rm ar}-\mathcal B_k^{(q+1)}}{q^2}
       &\longrightarrow-\Psi(1/4),\\
 \frac{\mathcal B_k^{\rm ar}}{q^2}&\longrightarrow C_B>0.
 \end{aligned}}
 \tag{KPR16}
\]
The scalar cancellation in the last line has the exact original
finite identity
$\mathcal B_k^{\rm ar}=\mathcal B_k^{(q+1)}
 -\mathfrak T_{k;q/4}+\delta_k^\sigma$ underneath it.
The mixed kernel and boundary receive the same equality through
IKO8 and BSL20, with every original observation direction and
off-diagonal Gram block retained. Thus the comparable-order
construction has produced and evaluated a positive $q^2$ return,
and its simultaneous baseline change is explicitly accounted for.
